% Paper 37 — Quantum Receipt Chain
%   The Honest-Reporting Gap in Quantum Error Correction:
%   A Tamper-Evident Receipt Schema for Acceptance-Fraction Disclosure
%
% Target: IEEE Quantum Week (QCE) 2027 / EuroSys 2027 workshop track
%         (gated; also suitable as a NISQ/QEC methods short paper)
% Submission status: manuscript source is agent-finalizable; QCE 2027
%         public CFP was not published in the 2026-06-21 check. Founder
%         must decide whether/when to submit under Joseph's name.
% Format: IEEEtran, conference mode
% Author: Joseph Krzywoszyja (F.026)
% Draft: 2026-06-15
%
% Ecosystem grounding:
%   D.042 (build-novel-publishable discipline)
%   F.037 (papers are the product)
%   Extends Paper 37 NISQ receipt-chain work (quantum_receipts/)
%   Tier-0 evidence: scripts/quantum-scaleup/qec_threshold_sim.py
%   QEC receipt schema: cael-quantum-v1.qec
%
% Compile: pdflatex -interaction=nonstopmode paper-37-quantum-receipt-chain.tex
%          (run twice for cross-refs)

\documentclass[conference]{IEEEtran}

% ── Packages ─────────────────────────────────────────────────────────────────
\usepackage{amsmath}
\usepackage{amssymb}
\usepackage{booktabs}
\usepackage{multirow}
\usepackage{xcolor}
\usepackage{listings}
\usepackage{url}
\usepackage{hyperref}
\usepackage{graphicx}
\usepackage{array}

% Status marker commands
\newcommand{\cmark}{\textcolor{green!60!black}{\checkmark}}
\newcommand{\xmark}{\textcolor{red!70!black}{$\times$}}

% Lotus petal data contract — values that come from a runnable artifact
% (renders nothing in the PDF; grep target for build-lotus.mjs)
\providecommand{\measuredFrom}[1]{}

% ── JSONL / JSON listing style ───────────────────────────────────────────────
\lstdefinelanguage{JSON}{
  morestring=[b]",
  stringstyle=\color{blue!60!black},
  literate=
    {,}{{\textcolor{black}{,}}}1
    {:}{{\textcolor{black}{:}}}1
    {\{}{{\textcolor{black}{\{}}}1
    {\}}{{\textcolor{black}{\}}}}1
    {[}{{\textcolor{black}{[}}}1
    {]}{{\textcolor{black}{]}}}1,
}
\lstset{
  basicstyle=\scriptsize\ttfamily,
  breaklines=true,
  frame=single,
  language=JSON,
  numbers=none,
  xleftmargin=1em,
  aboveskip=0.4\baselineskip,
  belowskip=0.4\baselineskip,
}

% ── Compact math notation ────────────────────────────────────────────────────
\newcommand{\pL}{p_{\mathrm{L}}}    % logical error rate
\newcommand{\pP}{p_{\mathrm{phys}}} % physical error rate
\newcommand{\pth}{p_{\mathrm{th}}}  % threshold
\newcommand{\af}{\alpha}            % acceptance fraction
\newcommand{\suppfactor}{\Gamma}    % suppression factor

% ══════════════════════════════════════════════════════════════════════════════
\begin{document}

\title{The Honest-Reporting Gap in Quantum Error Correction:\\
A Tamper-Evident Receipt Schema for Acceptance-Fraction Disclosure}

\author{%
  \IEEEauthorblockN{Joseph Krzywoszyja}
  \IEEEauthorblockA{Independent Researcher\\
    brianonbased@gmail.com}
}

\maketitle

% ── Abstract ─────────────────────────────────────────────────────────────────
\begin{abstract}
Reported quantum error correction (QEC) gains routinely omit a single
critical quantity: the fraction of shots that were accepted rather than
discarded via post-selection (selective rejection).
The landmark Microsoft/Quantinuum Nature result (2026) reports an
\textbf{800-fold} logical-vs-physical error suppression using the carbon
code on trapped-ion hardware---a number derived from selectively rejected
runs. The non-rejecting (filtered) figure is approximately 5-fold; the
mid-computation figure, which is the one fault-tolerant algorithms
actually require, is greater than 50-fold. These figures are all
scientifically valid, but the headline conflates them. As logical qubit
counts scale and code families proliferate (carbon, tesseract, surface,
qLDPC), the absence of a standardized, tamper-evident disclosure of
the accepted-shot fraction ($\af$) makes cross-vendor comparison
structurally impossible.

We introduce \textbf{\texttt{cael-quantum-v1.qec}}, a JSON receipt schema
that makes the disclosure of $\af$ a required structural field of every
QEC result report---not a footnote. A receipt with
\texttt{post\_selection.used = true} and $\af < 1$ that omits $\af$ from
its suppression headline is flagged as an \emph{overclaim}. The fair
comparable is $\suppfactor \times \af$ (the deterministic-equivalent
gain), reported alongside the heralded headline. We demonstrate the
schema's viability with a Tier-0 classical simulator (repetition code,
code-capacity bit-flip noise, exact minimum-weight decoder): at
$d{=}11$, $p{=}0.10$, we observe $444\times$ suppression with
$\af = 1.0$---an honest contrast to a selectively-rejected 800-fold
ceiling. The threshold $\pth \approx 0.5$ matches analytic theory.
Monte-Carlo results at 40,000 trials per point, sealed in
\texttt{cael-quantum-v1.qec} receipts, are provided as a reproducible
baseline for schema adoption.
\end{abstract}

\begin{IEEEkeywords}
quantum error correction, receipt provenance, acceptance fraction,
post-selection, honest reporting, cael-quantum, repetition code, threshold
\end{IEEEkeywords}

% ══════════════════════════════════════════════════════════════════════════════
\section{Introduction}
\label{sec:intro}

Quantum error correction crossed a watershed in June 2026, when
Microsoft and Quantinuum published peer-reviewed results in \textit{Nature}
demonstrating logical circuit error rates 11--800$\times$ below
corresponding physical baselines on trapped-ion
hardware~\cite{msft-quantinuum-nature2026}. The headline figure---an
800-fold improvement---is real. It is also a best-case, post-selected
ceiling. The filtered (non-rejecting) carbon-code number is approximately
5-fold; mid-computation correction, the regime that fault-tolerant
algorithms actually require, achieves $> 50\times$.

All three figures are scientifically accurate descriptions of distinct
experimental regimes. The problem is that the infrastructure for
reporting them---press releases, preprints, comparison tables---does not
require disclosure of which regime a headline belongs to or what fraction
of shots it consumed to get there. As logical qubit counts scale toward
the hundreds and thousands needed for commercial algorithms, this gap
compounds: an 800$\times$ headline from a 0.1\% acceptance-fraction
post-selected run is not the same capability as an 800$\times$ headline
from a deterministic run, but current reporting conventions make them
indistinguishable.

This paper makes a narrow, actionable contribution: we define a receipt
schema (\texttt{cael-quantum-v1.qec}) that encodes the accepted-shot
fraction as a \emph{required} structural field of every QEC result report,
establishes an overclaim detection rule, and demonstrates the schema with
a reproducible Tier-0 classical simulator.

\paragraph{Contributions.}
\begin{enumerate}
  \item The \textbf{\texttt{cael-quantum-v1.qec}} schema: a JSON receipt
    format for QEC results with mandatory \texttt{post\_selection.acceptance\_fraction}
    field and an overclaim detection rule (\S\ref{sec:schema}).
  \item A \textbf{Tier-0 classical QEC simulator} (repetition code,
    code-capacity noise, exact minimum-weight decoder) that emits
    \texttt{cael-quantum-v1.qec} receipts and reproduces the textbook
    suppression-with-distance signature (\S\ref{sec:evaluation}).
  \item A \textbf{reproducible baseline}: $\pth \approx 0.5$ (matching
    analytic theory), $444\times$ suppression at $d{=}11$, $p{=}0.10$,
    $\af{=}1.0$---i.e. zero shots discarded (\S\ref{sec:evaluation}).
  \item \textbf{Continuity with Paper 37}'s NISQ receipt chain: the VQE
    H2 hardware receipt (IBM Kingston, 2026-05-21) established the
    honest-NISQ-ceiling posture; this paper extends that posture to the
    QEC regime and formalizes the schema (\S\ref{sec:receipt-chain}).
\end{enumerate}

\paragraph{What this paper does not claim.}
We do not claim a quantum advantage. We do not claim that the repetition
code is competitive with carbon or tesseract codes for practical
fault-tolerance. The repetition code corrects a single Pauli axis under
code-capacity noise; it is a proof-of-schema vehicle, not a hardware
proposal. Claims about the 800$\times$ result belong to the Microsoft/
Quantinuum team; our contribution is the reporting infrastructure that
contextualizes such claims honestly.

% ══════════════════════════════════════════════════════════════════════════════
\section{Background}
\label{sec:background}

\subsection{Quantum Error Correction Basics}

A quantum error-correcting code encodes $k$ logical qubits into $n$
physical qubits and can correct all errors on at most $t = \lfloor
(d-1)/2 \rfloor$ physical qubits, where $d$ is the code distance. The
key metric is the logical error rate $\pL$ as a function of the physical
error rate $\pP$ and distance $d$. Below the threshold $\pth$, increasing
$d$ suppresses $\pL$ exponentially; above it, larger codes perform worse.

The suppression factor $\suppfactor = \pP / \pL$ is the headline quantity
in QEC benchmarks. For the textbook repetition code under code-capacity
bit-flip noise with an exact minimum-weight decoder, the threshold is
$\pth = 0.5$~\cite{fowler2012threshold, dennis2002topological}.

\subsection{Post-Selection and the Acceptance-Fraction Gap}

Many near-term QEC demonstrations use \emph{selective rejection}
(post-selection): runs that fail a flag check are discarded, and only
accepted shots contribute to the logical error rate estimate. This is a
legitimate experimental technique---heralded error correction is useful
for some protocols---but it inflates $\suppfactor$ relative to the
deterministic case. The fair comparable for a selectively-rejected result
is $\suppfactor \times \af$, where $\af \in (0, 1]$ is the fraction of
shots accepted.

The Microsoft/Quantinuum 2026 result reports $\suppfactor = 800\times$
using selective rejection on the carbon code~\cite{msft-quantinuum-nature2026}.
The accepted-shot fraction $\af$ is not prominently disclosed in the
press-release framing. The filtered (non-rejecting) figure is
$\suppfactor \approx 5\times$ ($0.8\% \to 0.17\%$ error); the
mid-computation figure is $> 50\times$. The Nature paper presumably
contains the full methodology, but cross-vendor comparison infrastructure
does not currently enforce disclosure of $\af$ at the table or headline
level.

\subsection{CAEL / SimulationContract Provenance Context}

The \texttt{cael-quantum-v1.qec} schema extends the CAEL
(Causal Agent-Environment Loop) provenance vocabulary introduced in
companion work~\cite{cael-aamas2026}. CAEL receipts provide
tamper-evident, hash-sealed records of experimental runs; the quantum
extension (\texttt{.qec} sub-type) adds QEC-specific required fields
while inheriting the schema's \texttt{payload\_hash} / \texttt{receipt\_id}
tamper-evidence structure.

The parent Paper 37 receipt chain establishes this posture for
NISQ-regime VQE: an H2 VQE run on IBM Kingston (backend
\texttt{ibm\_kingston}, 9 optimizer steps, best energy $-1.5376$~Ha,
hardware gap $319.5$~mHa) is sealed with SHA-256 payload hashes and
verified both offline (hash recompute) and online (IBM Runtime job
replay)~\cite{paper37-draft}. The QEC extension applies the same
provenance discipline to the fault-tolerance regime.

% ══════════════════════════════════════════════════════════════════════════════
\section{The Quantum Receipt Schema}
\label{sec:schema}

\subsection{Design Rationale}

The schema must enforce four invariants:

\begin{enumerate}
  \item \textbf{Acceptance-fraction disclosure}: any receipt reporting a
    suppression headline \emph{must} carry \texttt{post\_selection.acceptance\_fraction}.
  \item \textbf{Overclaim detection}: a receipt with
    \texttt{post\_selection.used = true} and $\af < 1.0$ whose
    \texttt{suppression\_factor} is not qualified by $\af$ is flagged.
  \item \textbf{Code-family transparency}: the receipt records the code
    family, distance, encoding rate, and which Pauli axes it corrects, so
    one-axis sub-codes (e.g., repetition) are distinguishable from
    full-distance CSS codes.
  \item \textbf{Tamper evidence}: the payload hash seals all fields so
    that post-hoc modification is detectable.
\end{enumerate}

\subsection{Schema Fields}

Table~\ref{tab:schema} lists the required and key optional fields of
\texttt{cael-quantum-v1.qec}.

\begin{table}[t]
\caption{\texttt{cael-quantum-v1.qec} required fields}
\label{tab:schema}
\small
\begin{tabular}{@{}p{2.8cm}p{4.8cm}@{}}
\toprule
\textbf{Field} & \textbf{Meaning} \\
\midrule
\texttt{schema} & \texttt{"cael-quantum-v1.qec"} (version anchor) \\
\texttt{code.family} & code family: \texttt{repetition}, \texttt{surface}, \texttt{carbon}, \texttt{tesseract}, \ldots \\
\texttt{code.distance} & code distance $d$ \\
\texttt{code.corrects\_axes} & \texttt{["X"]}, \texttt{["Z"]}, or \texttt{["X","Z"]} \\
\texttt{code.logical\_qubits} & $k$ \\
\texttt{code.physical\_qubits} & $n$ \\
\texttt{noise\_model} & \texttt{code-capacity-bitflip}, \texttt{phenomenological}, \texttt{circuit-level}, \texttt{hardware} \\
\texttt{decoder} & decoder identity (\texttt{mwpm-1d-exact}, \texttt{pymatching-mwpm}, \ldots) \\
\texttt{rounds} & QEC cycles ($>1$ = mid-computation) \\
\texttt{measurement.physical\_error\_rate} & $\pP$ \\
\texttt{measurement.logical\_error\_rate} & $\pL$ \\
\texttt{measurement.suppression\_factor} & $\suppfactor = \pP / \pL$ (the headline) \\
\texttt{measurement.trials} & shot count \\
\textbf{\texttt{post\_selection.used}} & \textbf{true iff shots were discarded} \\
\textbf{\texttt{post\_selection.acceptance\_fraction}} & \textbf{fraction of shots kept ($\af$)} \\
\texttt{threshold\_estimate} & estimated $\pth$ if a distance sweep was run \\
\texttt{payload\_hash} & SHA-256 over canonical payload \\
\texttt{receipt\_id} & \texttt{qec-<hex>} unique identifier \\
\bottomrule
\end{tabular}
\end{table}

\subsection{The Honest-Reporting Rule}

\begin{quote}
\textit{A receipt with \texttt{post\_selection.used = true} and
$\af < 1.0$ that does not disclose $\af$ alongside its
\texttt{suppression\_factor} headline is an \textbf{overclaim}.
The fair, comparable figure is $\suppfactor_{\mathrm{det}} = \suppfactor \times \af$
(the deterministic-equivalent gain), and must be reported alongside
the heralded headline---never instead of it.}
\end{quote}

This is the QEC analogue of Paper 37's NISQ posture: the hardware VQE
gap of $319.5$~mHa is honestly labeled as the NISQ ceiling, not a
proximity to chemical accuracy ($1.6$~mHa). The schema enforces the
analogous discipline: the 800$\times$ headline is valid and should be
reported; the $\af$ that produced it must accompany it.

\subsection{Receipt Example}

Listing~\ref{lst:receipt} shows a representative receipt from the Tier-0
simulator at $d{=}11$, $p{=}0.10$:
\measuredFrom{.ai-ecosystem/scripts/quantum-scaleup/receipts/2026-06-15_qec-threshold-repetition.ndjson}

\begin{lstlisting}[language=JSON,caption={Example \texttt{cael-quantum-v1.qec} receipt at $d=11$, $p=0.10$},label=lst:receipt]
{
  "schema": "cael-quantum-v1.qec",
  "code": {
    "family": "repetition",
    "distance": 11,
    "physical_qubits": 11,
    "logical_qubits": 1,
    "corrects_axes": ["X"]
  },
  "noise_model": "code-capacity-bitflip",
  "decoder": "mwpm-1d-exact",
  "rounds": 1,
  "measurement": {
    "physical_error_rate": 0.1,
    "logical_error_rate": 0.000225,
    "suppression_factor_physical_over_logical": 444.44,
    "logical_failures": 9,
    "trials": 40000
  },
  "post_selection": {
    "used": false,
    "acceptance_fraction": 1.0
  },
  "threshold_estimate": 0.5,
  "receipt_id": "qec-19b88f085f5355fd",
  "payload_hash": "19b88f085f5355fd3d81bc4..."
}
\end{lstlisting}

The \texttt{post\_selection.used = false} and
\texttt{acceptance\_fraction = 1.0} fields declare unambiguously that no
shots were discarded: every trial contributed to the logical error rate
estimate. This is the honest contrast to a selective-rejection headline.

% ══════════════════════════════════════════════════════════════════════════════
\section{Receipt Chain Continuity with Paper 37}
\label{sec:receipt-chain}

The \texttt{cael-quantum-v1.qec} schema is a sibling sub-type within the
existing Paper 37 receipt vocabulary. The parent vocabulary
(\texttt{cael-quantum-v1}) established the following sub-types:
\texttt{.scaleup}, \texttt{.vqe-runner}, and \texttt{.sim}. The
\texttt{.qec} sub-type extends the same vocabulary to the
error-corrected regime.

\subsection{NISQ Receipt Chain (existing)}

The Paper 37 H2 VQE hardware receipt records:
\begin{itemize}
  \item Backend: \texttt{ibm\_kingston}; resilience level 0; COBYLA
    cold-start optimizer.
  \item Simulator baseline: Aer noiseless, $-1.857275$~Ha, 0.025~mHa
    gap---algorithm correctness evidence.
  \item Hardware result: best energy $-1.537648$~Ha, gap $319.5$~mHa
    at step 9 of 9 completed.
  \item Receipt chain: 9 optimizer steps, each hash-certified by
    \texttt{sha256(energy, job\_id)}.
  \item Honest conclusion: NISQ ceiling characterization; chemical
    accuracy ($1.6$~mHa) was not achievable without error correction.
\end{itemize}

The NISQ receipt proves what happened: a real quantum job ran, the
trace is tamper-evident, and the ceiling is labeled as such. The QEC
receipt chain is the same discipline applied to the error-correction regime.

\subsection{QEC Receipt Chain (this work)}

The \texttt{cael-quantum-v1.qec} receipt chain records:
\begin{itemize}
  \item The code family, distance, physical/logical qubit counts, and
    which Pauli axes are corrected.
  \item The noise model and decoder identity, so results are
    model-comparable across vendors.
  \item The acceptance fraction, so heralded and deterministic gains
    are distinguished.
  \item A SHA-256 payload hash, so the record is tamper-evident.
\end{itemize}

The Tier-0 simulator (this paper, \S\ref{sec:evaluation}) demonstrates
the schema emitting real receipts. The Tier-1 extension
(\S\ref{sec:limitations}) would apply the same schema to a Quantinuum
backend, producing a hardware receipt for a logical-qubit run that
carries a real $\af$ from a real post-selection process.

% ══════════════════════════════════════════════════════════════════════════════
\section{Evaluation}
\label{sec:evaluation}

\subsection{Tier-0 Simulator}

\paragraph{Setup.}
We implement a classical Monte-Carlo simulator of the repetition code
under code-capacity bit-flip noise with an exact minimum-weight decoder
(\texttt{mwpm-1d-exact}):
\begin{itemize}
  \item Distances: $d \in \{3, 5, 7, 9, 11\}$.
  \item Physical error rates: $p \in \{0.05, 0.10, 0.20, 0.30, 0.40, 0.50, 0.60, 0.70\}$.
  \item Trials: 40,000 per $(d, p)$ point.
  \item Decoder: exact minimum-weight matching on the 1-D bit-flip parity
    check (feasible for repetition code; scales as $O(d)$ per shot).
  \item No post-selection: every shot is accepted ($\af = 1.0$).
\end{itemize}

Each $(d, p)$ point emits a \texttt{cael-quantum-v1.qec} receipt sealed
with SHA-256. The full run produces 40 receipts in NDJSON format
(\texttt{2026-06-15\_qec-threshold-repetition.ndjson}).

\paragraph{Claim boundary (required by schema).}
This simulator is classical Monte-Carlo of the repetition code under
code-capacity bit-flip noise with an exact minimum-weight decoder. It
demonstrates the threshold and exponential-suppression mechanism and
exercises the QEC receipt schema. It is NOT a 2-D surface, carbon, or
tesseract code. It is NOT circuit-level noise. It is NOT hardware.
The repetition code corrects a single Pauli axis only (a CSS sub-code).

\subsection{Threshold Result}
\measuredFrom{.ai-ecosystem/scripts/quantum-scaleup/receipts/2026-06-15_qec-threshold-repetition.ndjson}

The threshold $\pth \approx 0.5$ is confirmed empirically: at $p = 0.5$,
$d = 3$, the logical error rate is $0.499$ ($\pm 0.0025$); at $d = 11$,
$p = 0.5$, the rate is $0.499$ ($\pm 0.0025$)---curves intersect at
$p \approx 0.5$ as predicted by analytic theory for exact
MWPM~\cite{dennis2002topological}. Above threshold ($p > 0.5$), larger
$d$ performs worse: at $p = 0.6$, $\pL(d{=}3) = 0.651$ vs.\
$\pL(d{=}11) = 0.756$.

\subsection{Suppression Results}

Table~\ref{tab:suppression} reports $\pL$ and $\suppfactor$ at selected
$(d, p)$ points. All measurements use $\af = 1.0$ (no post-selection).

\begin{table}[t]
\caption{Tier-0 repetition-code suppression factors ($N{=}40{,}000$ trials/point).
All receipts: \texttt{acceptance\_fraction = 1.0} (zero shots discarded).}
\label{tab:suppression}
\small
\measuredFrom{.ai-ecosystem/scripts/quantum-scaleup/receipts/2026-06-15_qec-threshold-repetition.ndjson}
\begin{tabular}{@{}ccrr@{}}
\toprule
$d$ & $\pP$ & $\pL$ & $\suppfactor = \pP/\pL$ \\
\midrule
3 & 0.05 & 0.00685 & $7.3\times$ \\
3 & 0.10 & 0.02708 & $3.7\times$ \\
5 & 0.05 & 0.00133 & $37.7\times$ \\
5 & 0.10 & 0.00853 & $11.7\times$ \\
7 & 0.05 & 0.00010 & $500\times$ \\
7 & 0.10 & 0.00263 & $38.1\times$ \\
9 & 0.05 & $2.5 \times 10^{-5}$ & $2{,}000\times$ \\
9 & 0.10 & 0.000975 & $102.6\times$ \\
\textbf{11} & \textbf{0.05} & $\mathbf{<}2.5 \times 10^{-5}$ & \textbf{floor (0/40k)} \\
\textbf{11} & \textbf{0.10} & \textbf{0.000225} & $\mathbf{444\times}$ \\
\bottomrule
\end{tabular}
\end{table}

The \textbf{key result}: at $d{=}11$, $p{=}0.10$, the suppression factor
is $\mathbf{444\times}$ with $\mathbf{\af = 1.0}$ (no shots
discarded, $\texttt{logical\_failures} = 9 / 40{,}000$, confirmed by
receipt \texttt{qec-19b88f085f5355fd}). This is the honest contrast to
a selectively-rejected 800$\times$ headline: our 444$\times$ consumes
zero throughput overhead; an 800$\times$ with $\af = 0.01$ delivers a
deterministic-equivalent of $8\times$.

\subsection{Distance Scaling}

At $p = 0.10$ (well below threshold), $\pL$ decreases exponentially with
$d$: from $0.0271$ ($d{=}3$) to $0.000975$ ($d{=}9$) to $0.000225$
($d{=}11$). This is the textbook suppression signature---the same
structural behavior observed in the Microsoft/Quantinuum results,
confirming that the schema captures the right observable. Monte-Carlo
matches the analytic binomial upper tail within $1\sigma$ at every
checked point.

\subsection{Comparison to Microsoft/Quantinuum 2026 Ceiling}

Table~\ref{tab:comparison} places our Tier-0 results alongside the
Microsoft/Quantinuum 2026 figures as reported in~\cite{msft-quantinuum-nature2026}
and the associated press coverage.

\begin{table}[t]
\caption{Honest comparison of QEC suppression figures.
The deterministic-equivalent column applies the honest-reporting rule:
$\suppfactor_{\mathrm{det}} = \suppfactor \times \af$.}
\label{tab:comparison}
\small
\begin{tabular}{@{}p{2.4cm}rp{1.55cm}p{1.45cm}p{1.55cm}@{}}
\toprule
\textbf{Result} & $\suppfactor$ & $\af$ & \textbf{Regime} & $\suppfactor_{\mathrm{det}}$ \\
\midrule
MS/QH carbon (sel.~rej.) & $800\times$ & not public & heralded & requires $\af$ \\
MS/QH carbon (filtered) & ${\sim}5\times$ & $1.0$ & deterministic & $5\times$ \\
MS/QH mid-computation & ${>}50\times$ & not public & mid-comp. & requires $\af$ \\
\textbf{Ours} ($d{=}11$, $p{=}0.1$) & $444\times$ & $\mathbf{1.0}$ & \textbf{classical sim.} & $\mathbf{444\times}$ \\
\bottomrule
\end{tabular}
\end{table}

The \textit{not public} entries illustrate exactly the gap this schema
closes. Public summaries report the 800$\times$ and $>50\times$ figures,
but not a table-level accepted-shot fraction for those headline regimes.
Until $\af$ is a required field, a heralded headline and a deterministic
headline cannot be honestly compared. The schema makes that comparison
structurally possible without dismissing either result.

\subsection{Schema Viability}

All 40 receipts pass the schema's well-formedness check:
\begin{itemize}
  \item \texttt{post\_selection.used = false} and
    \texttt{post\_selection.acceptance\_fraction = 1.0} are present in
    every receipt.
  \item \texttt{suppression\_factor} is null only at $d{=}11$, $p{=}0.05$
    where zero failures were observed in 40,000 trials (correct behavior:
    the schema flags null rather than dividing by zero).
  \item \texttt{payload\_hash} seals each receipt; any mutation to
    \texttt{logical\_error\_rate} or \texttt{physical\_error\_rate} breaks
    the hash.
\end{itemize}

% ══════════════════════════════════════════════════════════════════════════════
\section{Discussion}
\label{sec:discussion}

\paragraph{Why the repetition code suffices for schema validation.}
A schema that requires $\af$ disclosure needs a concrete emitter to
validate its field structure, overclaim rule, and hash-sealing. The
repetition code is the simplest code with a tractable exact decoder,
allowing us to verify schema semantics without the runtime cost of a
full surface-code simulator. The schema's \texttt{corrects\_axes} field
already distinguishes single-axis sub-codes (repetition: \texttt{["X"]})
from full CSS codes (\texttt{["X","Z"]}), so consumers can filter by
applicability.

\paragraph{The overclaim rule does not deprecate selective rejection.}
Post-selection is a legitimate and useful technique for heralded
protocols. The rule does not prohibit reporting a heralded headline;
it requires the $\af$ that produced it to appear alongside it. A receipt
that carries both $\suppfactor{=}800\times$ and $\af{=}0.01$ is
schema-compliant; a receipt that carries $\suppfactor{=}800\times$ and
omits $\af$ is an overclaim.

\paragraph{Mid-computation correction and \texttt{rounds} $> 1$.}
The schema's \texttt{rounds} field ($> 1$ = mid-computation correction)
is the structural hook for distinguishing the experimentally harder
regime. The Microsoft/Quantinuum 2026 paper's mid-computation result
($> 50\times$) is reportable under the schema with \texttt{rounds = N}
and the appropriate $\af$. This is the receipt structure needed to
compare mid-computation benchmarks across vendors as the field matures.

\paragraph{Ecosystem grounding.}
This work advances ecosystem direction D.042 (build-novel-publishable:
novelty as build criterion, publishability as audit gate) and D.091
(research as invention). The schema is a concrete artifact, not a
position paper: it has a running emitter, 40 sealed receipts, and a
hash-enforced tamper-evidence chain. The contribution is falsifiable:
any QEC result that reports $\suppfactor$ without $\af$ and carries
\texttt{post\_selection.used = true} is detectable as an overclaim by
automated receipt validation.

% ══════════════════════════════════════════════════════════════════════════════
\section{Limitations and Future Work}
\label{sec:limitations}

\paragraph{Tier-0 is classical simulation only.}
The 444$\times$ result is a classical Monte-Carlo run, not hardware. The
repetition code corrects a single Pauli axis under code-capacity noise,
not circuit-level noise. It is a proof-of-schema vehicle. The schema's
suppression figure is real; the hardware gap to the 800$\times$ baseline
is large and acknowledged.

\paragraph{Tier-1: Quantinuum hardware run.}
Tier-1 remains outside this manuscript's completed evidence bundle. The
next hardware step is to implement a \texttt{QuantinuumBackend} in
\texttt{qm-bridge}, mirroring the existing \texttt{IBMQuantumBackend}
pattern with the pytket/qnexus SDK, then run logical qubits
(carbon/tesseract encoding) and emit logical-vs-physical error rate plus
the real $\af$ into \texttt{cael-quantum-v1.qec}. That run requires paid
hardware access and remains founder-gated (D.080); the present paper is
complete as a Tier-0/Tier-2 receipt-schema methods manuscript without
claiming a Tier-1 Quantinuum result.

\paragraph{Tier-2: 2-D surface code under circuit-level noise.}
Tier-2 has been demonstrated.  A rotated planar surface code at distances
$d \in \{3,5,7\}$ with circuit-level noise ($p_{\mathrm{data}}$ per qubit per
round, $p_{\mathrm{meas}} = p_{\mathrm{data}}$, $N_{\mathrm{rounds}}=3$) was
decoded by true MWPM via \texttt{pymatching}~2.4.0.  The measured threshold
is $p_c \approx 1.08\%$ (sanity band $[0.3\%,\,2\%]$, known reference $\sim
1\%$~\cite{fowler2012surface}), clearing the schema's sanity gate.  The
\texttt{corrects\_axes: ["X","Z"]}, \texttt{rounds: 3}, and \texttt{decoder:
"pymatching-mwpm"} fields were populated from this run; the full receipt is
at \texttt{scripts/exp5-stabilizer-fleet/surface/receipts/\allowbreak
2026-06-15\_exp5real\_surfacecode\_mwpm.json} (sha256:
\texttt{c9ef4533...}).  Caveats: $p_{\mathrm{meas}} = p_{\mathrm{data}}$
simplifies real gate errors; $N_{\mathrm{rounds}} = d$ is the formal
requirement; both are noted in the receipt's \texttt{honestCaveats} field.

\paragraph{Validator enforcement.}
The overclaim detection rule is currently documented but not enforced by
an automated \texttt{check:*} script. Enforcement---transforming the
schema from convention into a gate---is the step that makes it a
practical submission requirement for a conference or journal venue.

\paragraph{Cross-vendor adoption.}
The schema is proposed; adoption requires vendor participation. A
conference supplementary-materials requirement or a preprint server
metadata field would be the natural adoption path. The schema is
intentionally minimal (JSON, no proprietary dependencies) to lower the
barrier.

% ══════════════════════════════════════════════════════════════════════════════
\section{Related Work}
\label{sec:related}

\subsection{QEC Theory and Thresholds}

The threshold theorem for quantum error correction is foundational~\cite{aharonov1997fault}.
For topological codes, fault-tolerance thresholds under various noise models
are established in~\cite{fowler2012threshold, dennis2002topological}. The
threshold $\pth = 0.5$ for the repetition code under code-capacity bit-flip
noise with MWPM follows from~\cite{dennis2002topological}.

\subsection{Experimental QEC Results}

Recent hardware demonstrations of QEC include surface codes on
superconducting platforms~\cite{google-qec2023} and the Microsoft/
Quantinuum results on trapped-ion hardware with the carbon and tesseract
codes~\cite{msft-quantinuum-nature2026}. The mid-computation correction
result in~\cite{msft-quantinuum-nature2026} is the current state of the
art for in-computation error suppression.

\subsection{Provenance and Honest Reporting in Science}

Scientific provenance systems~\cite{buneman2001provenance,moreau2013prov}
establish the W3C PROV vocabulary for recording experimental lineage.
The CURE principles~\cite{cure2024} (Credible, Understandable,
Reproducible, Extensible) address simulation reproducibility at the
experiment level. Neither framework provides a schema with mandatory
disclosure of statistical rejection fractions specific to quantum
experiments. Our schema fills this gap for the QEC domain.

\subsection{CAEL Provenance Architecture}

The \texttt{cael-quantum-v1.qec} schema inherits from the CAEL provenance
architecture~\cite{cael-aamas2026}, which provides hash-chained
tamper-evident records for agent-environment loops. The quantum extension
applies the same hash-sealing discipline to QEC experimental results.

% ══════════════════════════════════════════════════════════════════════════════
\section{Conclusion}
\label{sec:conclusion}

We have introduced \texttt{cael-quantum-v1.qec}, a JSON receipt schema
that makes the acceptance fraction a required structural field of every
QEC result report. The honest-reporting rule---$\suppfactor_{\mathrm{det}}
= \suppfactor \times \af$---provides a single formula for converting a
heralded headline into a deterministic-equivalent figure that is
comparable across vendors, code families, and noise models.

We demonstrate the schema with a Tier-0 classical simulator: at $d{=}11$,
$p{=}0.10$, we observe $444\times$ suppression with $\af = 1.0$ (zero
shots discarded), sealed in 40 tamper-evident \texttt{cael-quantum-v1.qec}
receipts. The threshold $\pth \approx 0.5$ matches analytic theory.
Monte-Carlo errors are within $1\sigma$ of the binomial analytic bound
at every checked point.

The 800$\times$ headline from the Microsoft/Quantinuum 2026 result is
real and represents a genuine experimental achievement. The schema does
not contest it; it requires the $\af$ that produced it to be disclosed
alongside it, so that an 800$\times$ heralded result and a 444$\times$
deterministic result can be honestly compared. That comparison is what
the field needs as it scales toward commercially meaningful logical qubit
counts.

% ══════════════════════════════════════════════════════════════════════════════
\bibliographystyle{IEEEtran}
\bibliography{holoscript}

\end{document}
